\zlabel{firstpagetocount}       % DO NOT REMOVE! Used for counting number of pages of main text

% Part Labels Explanation:
% Labels in the form \label{chapter:<part_name>} are used to categorize chapters into specific parts of the thesis structure.
% These labels are essential for calculating the ratio of text dedicated to different thesis sections (e.g., Introduction, Methodology, Results, Discussion, Summary).
% Allowed labels:
% - chapter:introduction -> Introduction (<10% of the text).
% - chapter:method       -> Methodology (<20% of the text).
% - chapter:results      -> Results (30–40% of the text).
% - chapter:discussion   -> Analysis, Discussion, and Conclusions (30–40% of the text).
% - chapter:summary      -> Summary (<0.5 pages).
% How to use: Add a \label{chapter:<part_name>} for each chapter to indicate its part. 
% These labels ensure consistency and allow automated tests to validate the thesis structure.

\chapter{Introduction}\label{chapter:introduction} % This command creates a numbered chapter titled "Sissejuhatus" (Estonian for "Introduction"). The \label{chapter:introduction} assigns a label to the chapter, allowing you to reference it later (e.g., \ref{chapter:introduction}).
Unmanned aerial vehicles (UAVs) have become very popular tools across military, commercial, and research sectors.\cite{drones} Accurate aerodynamic analysis is critical for optimizing their performance, particularly when evaluating the impact of external payloads on drag and efficiency. Computational Fluid Dynamics (CFD) offers a cost-effective alternative to wind tunnel testing for aerodynamic predictions. However, CFD results must be properly validated against experimental data before they can be reliably applied to new design configurations.

In aerospace engineering, commercial CFD programmes such as ANSYS Fluent dominate due to their convenient user interface and comprehensive support. OpenFOAM, an open-source CFD toolbox, provides a compelling alternative with zero costs and complete source code access. Despite its capabilities, OpenFOAM remains underutilized in aerospace applications compared to commercial software, partly due to difficult learning process and concerns about validation. 

This thesis addresses this gap by validating OpenFOAM against both wind tunnel experimental data and ANSYS Fluent predictions for a small-scale delta wing aircraft. The validation approach uses published wind tunnel measurements of a delta wing configuration with dimensions comparable to the target drone developed by CAFA Tech. By comparing aerodynamic coefficients from OpenFOAM simulations against both experimental results and ANSYS Fluent computations, the accuracy of the open-source solver can be quantitatively assessed. Delta wing configurations are particularly suitable for validation due to their complex vortical flow structures that challenge CFD accuracy.

Once the OpenFOAM simulation framework is validated, the same computational setup will be applied to analyze the CAFA Tech fixed-wing drone in two configurations: the baseline clean aircraft and the configuration with two rocket-shaped external stores mounted beneath the fuselage. The primary objective is to evaluate a drag increase by the external stores. This information is essential for CAFA Tech's high speed drone development, because external payloads reduce aerodynamic efficiency. 
 %  This command inserts the content from the file introduction.tex located in the chapters folder.

\chapter{CFD Fundamentals for Aerodynamic Simulations}\label{chapter:method}
In this chapter, theoretical basis for the OpenFOAM simulation is described. 
\section{Reynolds-Averaged Navier-Stokes (RANS) Equations}
\label{sec:rans_theory}

The motion of any fluid is governed by the Navier--Stokes equations, which describe the conservation of mass and momentum. For an incompressible fluid with constant density, the continuity equation requires that the velocity field is divergence-free, and the momentum equation balances the rate of change of velocity with pressure gradients, viscous forces, and any external forces \cite{Pope2000TurbulentFlows,VersteegMalalasekera2007CFD}. These equations are exact and apply to all flow regimes, including turbulent flows. However, solving them directly for turbulent flows requires resolving all small-scale fluctuations in velocity and pressure, which demands extremely fine meshes and very small time steps. For most engineering applications this is computationally too expensive \cite{Pope2000TurbulentFlows}.

To make the problem practical, the Reynolds-Averaged Navier--Stokes (RANS) approach applies a time-averaging procedure using Reynolds decomposition \cite{Pope2000TurbulentFlows,Wilcox2006TurbulenceModeling}. Each flow variable is split into a mean (time-averaged) component and a fluctuating component. For the velocity field this is written as
\begin{equation}
u_i = \overline{u}_i + u'_i ,
\end{equation}
where $\overline{u}_i$ is the mean velocity and $u'_i$ is the turbulent fluctuation. The same decomposition is applied to pressure. Substituting these decomposed variables into the Navier--Stokes equations and applying time-averaging yields the RANS equations \cite{Pope2000TurbulentFlows,VersteegMalalasekera2007CFD}.

For incompressible flow, the averaged continuity equation retains the same form as the instantaneous equation:
\begin{equation}
\frac{\partial \overline{u}_i}{\partial x_i} = 0 .
\end{equation}
The averaged momentum equation becomes:
\begin{equation}
\rho\left(
\frac{\partial \overline{u}_i}{\partial t}
+
\overline{u}_j \frac{\partial \overline{u}_i}{\partial x_j}
\right)
=
-\frac{\partial \overline{p}}{\partial x_i}
+
\mu \frac{\partial^2 \overline{u}_i}{\partial x_j \partial x_j}
-
\rho \frac{\partial \left(\overline{u'_i u'_j}\right)}{\partial x_j} .
\end{equation}

This equation is similar to the laminar Navier--Stokes momentum equation, except for the additional term involving $\overline{u'_i u'_j}$ \cite{Pope2000TurbulentFlows,VersteegMalalasekera2007CFD}. The quantity $-\rho \overline{u'_i u'_j}$ is known as the Reynolds stress tensor. It represents the influence of turbulent fluctuations on the mean flow and appears due to averaging of the nonlinear convective term. The Reynolds stress tensor introduces six additional unknowns (it is symmetric), but the RANS equations provide no extra equations to solve for them. This is known as the turbulence closure problem \cite{Pope2000TurbulentFlows,Wilcox2006TurbulenceModeling}.

To close the system, a turbulence model is required. A common engineering approach is the eddy-viscosity hypothesis, originally proposed by Boussinesq, which assumes that the Reynolds stresses can be related to the mean rate of strain through a turbulent (eddy) viscosity $\mu_t$ \cite{VersteegMalalasekera2007CFD,Wilcox2006TurbulenceModeling}. With this assumption, the momentum equation can be written in a form similar to the laminar case, but using an effective viscosity that includes both molecular and turbulent contributions ($\mu+\mu_t$). The turbulence model therefore provides $\mu_t$ (or equivalently $\nu_t$) throughout the flow field.

This is the approach used in OpenFOAM. The \texttt{simpleFoam} solver solves the steady-state, incompressible RANS equations using the SIMPLE pressure--velocity coupling algorithm \cite{VersteegMalalasekera2007CFD}. The mean velocity and pressure fields are computed, while the turbulent viscosity is obtained from the selected turbulence model. When the Spalart--Allmaras model is used, a single transport equation for a modified turbulent viscosity $\tilde{\nu}$ is solved and $\nu_t$ is computed from $\tilde{\nu}$ \cite{SpalartAllmaras1992SA,SpalartAllmaras1994SAJournal}. When the $k$--$\omega$ SST model is used, two transport equations are solved for the turbulent kinetic energy $k$ and the specific dissipation rate $\omega$, and $\nu_t$ is computed from $k$ and $\omega$ \cite{Menter1994SST}. In both cases, the underlying mean-flow equations solved by \texttt{simpleFoam} remain the same; only the method of evaluating the turbulent viscosity changes.
\section{Turbulence Modeling}
\label{sec:turbulence_modeling}

As described in Section~\ref{sec:rans_theory}, the RANS equations introduce the Reynolds stress tensor, which adds unknown quantities that must be modelled. The most widely used approach in engineering CFD is based on the Boussinesq eddy-viscosity hypothesis, which assumes that the turbulent stresses can be represented using a single scalar quantity, the turbulent viscosity $\mu_t$ \cite{Pope2000TurbulentFlows,VersteegMalalasekera2007CFD,Wilcox2006TurbulenceModeling}. The role of a turbulence model is therefore to calculate $\mu_t$ at every point in the flow field. Different models achieve this by solving different sets of transport equations for turbulence-related variables.

Turbulence models are commonly classified by the number of additional transport equations they solve. Zero-equation (algebraic) models compute $\mu_t$ directly from the mean flow without solving additional equations. One-equation models solve a single transport equation for a turbulence variable. Two-equation models solve two transport equations, which allows the turbulence velocity scale and turbulence length scale to be determined independently \cite{VersteegMalalasekera2007CFD,Wilcox2006TurbulenceModeling}. In general, models with more equations can represent a wider range of flow situations, but they also introduce more empirical constants and require additional computational effort.

For this thesis, two turbulence models were selected: the one-equation Spalart--Allmaras (SA) model and the two-equation $k$--$\omega$ Shear Stress Transport (SST) model. Both models are widely used in aerospace CFD and were also used in the reference ANSYS Fluent simulations by Giannelis et al.~\cite{Giannelis2023SSAMGen5}. Using the same models in OpenFOAM supports a direct comparison of turbulence-closure behaviour between the solvers.

\subsection{$k$--$\omega$ SST Model}
\label{sec:sst_model}

The $k$--$\omega$ SST model was developed by Menter \cite{Menter1994SST} to address limitations of two earlier two-equation models: the standard $k$--$\omega$ model and the standard $k$--$\varepsilon$ model. The $k$--$\omega$ model performs well in the near-wall region and within boundary layers, but it is sensitive to the freestream value of $\omega$, meaning that results may depend on the inlet boundary condition for $\omega$ \cite{Wilcox2006TurbulenceModeling,Menter1994SST}. In contrast, the $k$--$\varepsilon$ model is less sensitive to freestream conditions, but it typically performs worse near walls and may give poor predictions for flows with adverse pressure gradients and separation \cite{VersteegMalalasekera2007CFD,Wilcox2006TurbulenceModeling}.

The SST model combines the advantages of both approaches using a blending function. Near walls the model behaves like a $k$--$\omega$ formulation, and it gradually transitions to a $k$--$\varepsilon$-type behaviour (rewritten in terms of $\omega$) in the outer boundary layer and in the freestream \cite{Menter1994SST}. In addition to the blending, the SST formulation introduces a modified turbulent-viscosity definition that accounts for transport of turbulent shear stress. This improves predictions for flows with adverse pressure gradients and separation, which are important in aerodynamic applications \cite{Menter1994SST}.

The SST model solves two transport equations: one for turbulent kinetic energy $k$, which represents the energy content of turbulent fluctuations, and one for the specific dissipation rate $\omega$, which controls the rate at which turbulent energy is dissipated \cite{VersteegMalalasekera2007CFD,Menter1994SST}. The turbulent viscosity is then computed in the form
\begin{equation}
\mu_t = \rho \, \frac{k}{\omega},
\end{equation}
with an additional SST limiter applied through the shear-stress transport modification to prevent unrealistically large $\mu_t$ in regions of high production (for example near stagnation points or in strong vortex cores) \cite{Menter1994SST}.

In OpenFOAM, the $k$--$\omega$ SST model is available as a standard selection in \texttt{turbulenceProperties}. When selected, \texttt{simpleFoam} solves the additional transport equations for $k$ and $\omega$ alongside the RANS momentum and continuity equations. The case setup therefore requires the initial and boundary condition fields \texttt{k} and \texttt{omega} in the \texttt{0} directory. Near-wall treatment for $\omega$ is typically applied using \texttt{omegaWallFunction}, while $k$ commonly uses \texttt{kqRWallFunction} on solid surfaces \cite{VersteegMalalasekera2007CFD}.

\subsection{Spalart--Allmaras Model}
\label{sec:sa_model}

The Spalart--Allmaras (SA) model was developed specifically for aerodynamic applications \cite{SpalartAllmaras1992SA,SpalartAllmaras1994SAJournal}. Unlike two-equation models, it solves one transport equation for a modified turbulent-viscosity variable, usually denoted as $\tilde{\nu}$ (called \texttt{nuTilda} in OpenFOAM). The turbulent viscosity $\mu_t$ is then computed from $\tilde{\nu}$ using a damping function to ensure correct near-wall behaviour \cite{SpalartAllmaras1992SA,SpalartAllmaras1994SAJournal}.

The main motivation behind the SA model was to provide a simple and computationally efficient turbulence model suitable for external aerodynamic flows over wings and fuselages at different angles of attack \cite{SpalartAllmaras1992SA}. Since only one additional equation is solved, the model typically requires less computational effort per iteration than two-equation models and is often numerically stable, which helps convergence in complex geometries \cite{Wilcox2006TurbulenceModeling}.

The transport equation for $\tilde{\nu}$ includes production, diffusion, and destruction terms \cite{SpalartAllmaras1992SA,SpalartAllmaras1994SAJournal}. The production term depends on the local flow rotation (vorticity/strain), meaning that turbulent viscosity is generated in regions such as boundary layers and vortices. The destruction term ensures that $\tilde{\nu}$ is reduced in regions away from walls where turbulence production is weak. The balance between these terms determines the distribution of turbulent viscosity in the flow field \cite{Wilcox2006TurbulenceModeling}.

A known characteristic of the SA model is that it can overpredict turbulent viscosity inside strong vortex cores, which may increase numerical diffusion and weaken vortical structures \cite{Giannelis2023SSAMGen5}. Nevertheless, the reference study by Giannelis et al.~\cite{Giannelis2023SSAMGen5} reported that the SA model provided strong overall correlation with wind tunnel measurements for the SSAM-Gen5 geometry, particularly at higher angles of attack.

In OpenFOAM, the SA model is selected in \texttt{turbulenceProperties}. When active, \texttt{simpleFoam} solves the single transport equation for $\tilde{\nu}$ alongside the mean-flow equations. The case setup requires the turbulence field \texttt{nuTilda} in the \texttt{0} directory. At solid walls, \texttt{nuTilda} is typically set to zero, while the turbulent viscosity field \texttt{nut} uses a wall function boundary condition such as \texttt{nutUSpaldingWallFunction} to model near-wall behaviour \cite{VersteegMalalasekera2007CFD}.

Both the SA and SST models provide the turbulent viscosity that is added to the molecular viscosity in the RANS momentum equation. From the perspective of the \texttt{simpleFoam} solver, the mean-flow equations remain the same regardless of turbulence model choice; the difference is the method used to compute $\mu_t$ and the number of additional transport equations solved.


\section{Finite Volume Method and Discretization}

\label{sec:fvm_discretization}

The finite volume method (FVM) is the numerical approach used by OpenFOAM to solve the RANS equations. In this method, the computational domain is divided into a large number of small control volumes (cells), and the governing equations are integrated over each cell \cite{Patankar1980NumericalHeatTransfer,VersteegMalalasekera2007CFD}. This integration ensures that mass, momentum, and other conserved quantities are preserved at the discrete level, which is one of the main advantages of the finite volume approach \cite{VersteegMalalasekera2007CFD,FerzigerPericStreet2019CMFD}.

For each cell, the volume integrals of the governing equations are converted into surface integrals over the cell faces using the Gauss divergence theorem \cite{VersteegMalalasekera2007CFD}. The values of velocity, pressure, and turbulence variables at the cell faces must then be estimated from the values stored at the cell centres. This estimation is called discretization, and the chosen discretization schemes affect both the accuracy and numerical stability of the solution \cite{Patankar1980NumericalHeatTransfer,VersteegMalalasekera2007CFD}.

In the present simulations, second-order accurate schemes were used for the discretization of convective terms. Second-order schemes provide a practical balance between accuracy and stability for external aerodynamic flows \cite{VersteegMalalasekera2007CFD,FerzigerPericStreet2019CMFD}. For diffusive terms, second-order central differencing was applied, which is a standard choice in OpenFOAM. The gradients required for both convective and diffusive terms were computed using the Gauss linear method.

The coupling between pressure and velocity in incompressible flow was handled using the SIMPLE (Semi-Implicit Method for Pressure-Linked Equations) algorithm \cite{Patankar1980NumericalHeatTransfer,VersteegMalalasekera2007CFD}. In SIMPLE, the momentum equations are first solved using an estimated pressure field to obtain an intermediate velocity field. A pressure-correction equation is then solved to update the pressure and correct the velocity so that continuity is satisfied. This procedure is repeated iteratively until changes between successive iterations become sufficiently small and the solution is considered converged \cite{Patankar1980NumericalHeatTransfer,VersteegMalalasekera2007CFD}. In OpenFOAM, the SIMPLE algorithm is implemented in the \texttt{simpleFoam} solver, which was used for all steady-state simulations in this work. The specific numerical schemes and solver settings used in this thesis are provided in Appendix~[X].
\section{OpenFOAM Framework and Solvers}
\label{sec:openfoam_framework}

\subsection*{Overview}
OpenFOAM (Open Field Operation And Manipulation) is a free and open-source \texttt{C++} toolbox for computational fluid dynamics (CFD) and, more generally, computational continuum mechanics \cite{Weller1998FOAM,Jasak2009OpenFOAM}. The core design ideas of the FOAM framework were published by Weller et al. in 1998 \cite{Weller1998FOAM}. OpenFOAM was released as free open-source software under the GNU General Public License (GPL) in December 2004 \cite{OpenFOAMOpenSource}. Because it is distributed under the GPL, the software can be used, modified, and redistributed without license fees, and the full source code is available for inspection and development \cite{OpenFOAMLicensing,OpenFOAMOpenSource}.

\subsection*{Software architecture}
OpenFOAM is not a single graphical program. It is a collection of libraries and applications (utilities and solvers) that are typically executed from the command line \cite{Weller1998FOAM,Jasak2009OpenFOAM}. The framework is built using object-oriented programming, where mathematical objects such as vectors, tensors, and fields are represented as \texttt{C++} classes. This design allows solver code to closely follow the mathematical form of the equations being solved \cite{Weller1998FOAM,Jasak2009OpenFOAM}.

A typical OpenFOAM case is organised into three main directories \cite{OpenFOAMUserGuide}:
\begin{itemize}
    \item \texttt{0/} --- initial and boundary condition files for each field (e.g., velocity, pressure, turbulence variables),
    \item \texttt{constant/} --- mesh in \texttt{polyMesh/} and physical property files (e.g., viscosity and turbulence model selection),
    \item \texttt{system/} --- numerical schemes, solver settings, and run control parameters.
\end{itemize}
The case setup is performed using plain-text dictionary files. For post-processing and visualisation, OpenFOAM results can be opened directly in ParaView \cite{Ahrens2005ParaView,OpenFOAMUserGuide}.

\subsection*{Meshing utilities}
OpenFOAM includes built-in utilities for mesh generation. A common workflow is to create a simple background mesh using \texttt{blockMesh} and then refine and fit the mesh to an STL surface geometry using \texttt{snappyHexMesh} \cite{OpenFOAMUserGuide,OpenFOAMsnappyHexMeshGuide}. The \texttt{snappyHexMesh} process typically includes three stages: castellation (cell splitting/refinement), surface snapping, and boundary-layer addition \cite{OpenFOAMsnappyHexMeshGuide,OpenFOAMsnappyHexMeshV13}. The mesh generation procedure used in this thesis follows this standard approach (see Section~3.3).


\subsection*{Solver selection}
OpenFOAM provides many solvers, each designed for a specific class of problems (e.g., incompressible, compressible, multiphase, heat transfer) \cite{OpenFOAMUserGuide}. For incompressible flows, commonly used solvers include:
\begin{itemize}
    \item \texttt{simpleFoam} --- steady-state incompressible solver using the SIMPLE algorithm \cite{OpenFOAMsimpleFoamGuide},
    \item \texttt{pisoFoam} --- transient incompressible solver using the PISO algorithm \cite{OpenFOAMUserGuide},
    \item \texttt{pimpleFoam} --- transient incompressible solver using the PIMPLE algorithm (combination of PISO and SIMPLE) \cite{OpenFOAMUserGuide}.
\end{itemize}
For compressible flows, OpenFOAM provides solvers such as \texttt{rhoSimpleFoam} (steady-state) and \texttt{rhoPimpleFoam} (transient), which additionally solve an energy equation and account for density variations \cite{OpenFOAMUserGuide}.

\subsection*{Why \texttt{simpleFoam} was chosen}
In this thesis, \texttt{simpleFoam} was selected for all simulations for three reasons.

First, the simulated flow conditions correspond to low-speed incompressible flow. The freestream velocity is $20~\mathrm{m/s}$, which corresponds to a Mach number of approximately $M \approx 0.06$. At such low Mach numbers, density variations are negligible and an incompressible solver is appropriate \cite{VersteegMalalasekera2007CFD}.

Second, the study focuses on time-averaged aerodynamic coefficients (lift, drag, and pitching moment), rather than on the time evolution of the flow. A steady-state solver is therefore sufficient and computationally more efficient than a transient solver for this purpose. The \texttt{simpleFoam} solver iterates to a converged solution using SIMPLE \cite{OpenFOAMsimpleFoamGuide}.

Third, the reference ANSYS Fluent simulations by Giannelis et al.~\cite{Giannelis2023SSAMGen5} used a steady-state, pressure-based solver with SIMPLE-type coupling. Using \texttt{simpleFoam} therefore supports a direct comparison between the two CFD frameworks.

It is noted that steady-state RANS simulations have limitations at high angles of attack where the flow can become strongly separated and unsteady. In such cases the solution may not converge to a perfectly steady state and force coefficients can oscillate between iterations. For these cases, time-averaged coefficient values were reported by averaging over a stable interval of the final iterations, consistent with the approach described in the reference study \cite{Giannelis2023SSAMGen5}.

\subsection*{Comparison with ANSYS Fluent}
ANSYS Fluent is a commercial CFD package that provides a complete graphical workflow for geometry import, meshing, solver setup, and post-processing, supported by extensive documentation and commercial support \cite{AnsysFluentTheoryGuide}. OpenFOAM and ANSYS Fluent share the same core numerical approach: both solve the governing equations using the finite volume method and both provide SIMPLE-type pressure--velocity coupling for steady incompressible simulations \cite{VersteegMalalasekera2007CFD,OpenFOAMsimpleFoamGuide,AnsysFluentTheoryGuide}. Both codes also provide the turbulence models used in this thesis (Spalart--Allmaras and $k$--$\omega$ SST). Practical differences are mainly related to the user interface (text-based in OpenFOAM versus GUI-based in Fluent), meshing tools, and implementation details of the numerical methods. The validation results in Chapter~4 show that despite these differences, comparable aerodynamic predictions can be achieved for the same geometry and flow conditions.

\chapter{OpenFOAM Simulation Setup}\label{chapter:results}

\section{Validation Strategy and Reference Cases}
Before applying any CFD solver to a new or unfamiliar geometry, it is essential to first verify that the solver produces reliable results for a known case. Without this step, there is no way to judge whether the predictions for the actual target geometry are physically meaningful or affected by errors in the simulation setup. This is especially important when using OpenFOAM, because unlike commercial software such as ANSYS Fluent, it does not provide a graphical interface that guides the user through the setup process. Every setting must be configured manually through text files, which increases the risk of mistakes in boundary conditions, solver parameters, or mesh configuration. A structured validation study helps to identify and eliminate such errors before moving to the actual CAFA drone simulation.

The validation approach used in this thesis is based on a three-way comparison: OpenFOAM simulation results are compared against both wind tunnel experimental data and ANSYS Fluent predictions for the same geometry and flow conditions. This approach has two main advantages. First, comparison with experimental data shows how close the simulation is to real physical behaviour. Second, comparison with ANSYS Fluent results allows a direct evaluation of OpenFOAM against an established and widely trusted commercial solver. If OpenFOAM produces results that are close to both the experiment and ANSYS Fluent, it can be considered sufficiently accurate for the subsequent drone analysis.

The reference case selected for this validation is the SSAM-Gen5 model (Sydney Standard Aerodynamic Model -- Generation~5), a generic delta wing aircraft geometry. This model was chosen for several reasons. First, it is a delta wing configuration with sharp leading edges that produces complex vortex-dominated flow at moderate and high angles of attack \cite{Giannelis2023SSAMGen5}. This makes it a challenging test case for any CFD solver, because capturing vortex formation, interaction, and breakdown correctly requires good mesh quality and appropriate solver settings. If OpenFOAM can handle this level of flow complexity, it can be expected to perform well for the simpler flow around the CAFA Tech drone. Second, the SSAM-Gen5 model length is $0.75~\mathrm{m}$, which is comparable in scale to the CAFA Tech drone. This means that the Reynolds number range and mesh resolution requirements are similar, making the validation directly relevant to the final application. Third, all geometry files, experimental data, and computational meshes are publicly available through the Zenodo repository \cite{Giannelis2023SSAMGen5}, which ensures full transparency and reproducibility of the validation process.

The experimental reference data were obtained in subsonic wind tunnel tests at the University of Sydney at a freestream velocity of $20~\mathrm{m/s}$, corresponding to a Reynolds number of approximately $3 \times 10^{5}$ based on the mean aerodynamic chord \cite{Giannelis2023SSAMGen5}. Force and moment measurements were recorded across a range of angles of attack from $-5^{\circ}$ to $30^{\circ}$ \cite{Giannelis2023SSAMGen5}. The experimental results confirmed that the SSAM-Gen5 exhibits aerodynamic behaviour typical of slender delta wing aircraft, including monotonically increasing lift without abrupt stall and a pitching moment break near $12^{\circ}$ angle of attack \cite{Giannelis2023SSAMGen5}.

The ANSYS Fluent reference results were taken from the same publication by Giannelis et al.\ \cite{Giannelis2023SSAMGen5}. These simulations used the steady, incompressible, pressure-based solver with second-order spatial discretisation \cite{Giannelis2023SSAMGen5}. Several turbulence models were tested, including Spalart--Allmaras (SA), $k$--$\omega$ SST, and Realisable $k$--$\varepsilon$, both with and without curvature correction \cite{Giannelis2023SSAMGen5}. The results showed that the turbulence model choice has a strong effect on the predicted aerodynamic coefficients at high angles, with the SA model providing the best correlation with the experimental data overall \cite{Giannelis2023SSAMGen5}. For the validation in this thesis, the OpenFOAM results are compared against the ANSYS Fluent predictions obtained with both the $k$--$\omega$ SST and Spalart--Allmaras models, as these are the same two turbulence models used in the OpenFOAM simulations.

To ensure a fair comparison, all three data sources (wind tunnel, ANSYS Fluent, and OpenFOAM) use the same model geometry, the same freestream velocity of $20~\mathrm{m/s}$, and the same angle of attack range. The same reference area and reference length are used for computing the aerodynamic coefficients. The only differences lie in the mesh generation approach (\texttt{snappyHexMesh} in OpenFOAM versus Pointwise in ANSYS Fluent) and the internal numerical implementation of the solvers. These are precisely the differences that the validation study is designed to evaluate.

\section{Geometry: SSAM-Gen5 Delta-Wing Configuration}
In this section the geometry and experimental setup from Giannelis et al. \cite{Giannelis2023SSAMGen5} is summarized.
The geometry used for validation is the SSAM-Gen5 model, a generic aircraft designed to represent the key shape features of modern fifth-generation high-performance platforms. The overall shape is based on the general layout of the F-22 Raptor, but with significant simplifications to keep the geometry easy to manufacture and to model in CFD.

The full-scale vehicle length is $19~\mathrm{m}$, and the wind tunnel model is scaled to $0.75~\mathrm{m}$, which corresponds to approximately a 1:25 scale. The main reference values for the scaled model are summarised in Table~\ref{tab:ssam_ref_values}.

\begin{table}[H]
\centering
\caption{Reference values for the SSAM-Gen5 model at $0.75~\mathrm{m}$ scale \cite{Giannelis2023SSAMGen5}.}
\label{tab:ssam_ref_values}
\begin{tabular}{lcc}
\hline
\textbf{Parameter} & \textbf{Value} & \textbf{Unit} \\
\hline
Reference area          & 0.1091 & $\mathrm{m^2}$ \\
Mean aerodynamic chord  & 0.2265 & $\mathrm{m}$ \\
Span                    & 0.5350 & $\mathrm{m}$ \\
\hline
\end{tabular}
\end{table}

The wing profiles are four-percent-thick biconvex sections with sharp leading edges. No camber or spanwise twist is applied to the wings, which keeps the geometry as simple as possible while still producing realistic vortex-dominated flow. The sweep angle of the main delta wing and the horizontal stabiliser is approximately $42^{\circ}$, while the twin vertical stabilisers are swept back about $25^{\circ}$ and canted outwards at $27.5^{\circ}$. The aircraft includes engine intake openings, but these are simplified and sealed in the wind tunnel model. Thrust nozzles at the rear of the fuselage are retained in the geometry.

The sharp leading edges are an important feature of this geometry. Unlike wings with rounded leading edges, sharp edges force flow separation to occur at a fixed location along the entire leading edge, which leads to the formation of strong vortices over the upper surface of the wing. These vortices produce additional lift at moderate and high angles of attack, known as vortex lift. At the same time, the vortical flow structures make the aerodynamic predictions sensitive to mesh quality, turbulence model choice, and solver settings. This is what makes the SSAM-Gen5 a demanding and therefore useful validation case.

The wind tunnel model was fabricated from 3D-printed PLA parts, which were assembled, sanded, and coated with resin to provide a smooth surface finish. The model was designed with removable wings, fins, and horizontal stabilisers to allow testing of different configurations.

The STL geometry file used in the OpenFOAM simulations was downloaded directly from the publicly available Zenodo repository. This ensures that the exact same geometry is used in both the OpenFOAM and ANSYS Fluent simulations, eliminating any differences that could arise from geometry reconstruction or approximation.

\section{Computational Domain and Mesh Generation}
The computational domain for the simulations was defined as a rectangular box around the delta-wing model. The domain size was selected to reduce the influence of the outer boundaries on the flow around the wing, even at high angles of attack (above $20^\circ$). A sufficiently long downstream distance also helps to avoid reverse flow at the outlet and allows the wake behind the aircraft to develop inside the domain.

The delta-wing model length is $0.75~\mathrm{m}$. The domain limits were defined in Cartesian coordinates as:
\begin{itemize}
    \item $x$-direction: from $-3.0~\mathrm{m}$ to $7.5~\mathrm{m}$,
    \item $y$-direction: from $0.0~\mathrm{m}$ to $3.0~\mathrm{m}$,
    \item $z$-direction: from $-3.0~\mathrm{m}$ to $3.0~\mathrm{m}$.
\end{itemize}
Only half of the geometry was simulated by applying a symmetry plane at $y = 0$, which reduces the computational cost while remaining valid for a symmetric geometry and flow setup.

The background mesh was generated using the \texttt{blockMesh} utility, creating a three-dimensional hexahedral block. Boundary patches were assigned to each outer face of the domain, including \texttt{inlet} (front patch), \texttt{outlet} (rear patch), \texttt{lowerWall}, \texttt{upperWall}, and \texttt{sideWall}. A \texttt{symmetryPlane} boundary condition was defined to enable the half-domain simulation.

After the background mesh and boundaries were created, the \texttt{surfaceFeatureExtract} utility was used to extract feature edges from the delta-wing \texttt{.stl} geometry. This step helps preserve sharp geometric details (for example, leading and trailing edges) during the refinement and snapping stages. The main parameter of this utility is \texttt{includedAngle}, which identifies feature edges based on the angle between adjacent surface normals. For the delta-wing geometry, an \texttt{includedAngle} of $176^\circ$ was selected after several trials to capture sharp wing edges reliably.

Mesh generation was then performed using the \texttt{snappyHexMesh} utility, which produces a predominantly hexahedral mesh fitted to the surface geometry. The initial setup for this step was based on the official OpenFOAM \texttt{motorBike} tutorial and then adjusted through multiple attempts to obtain acceptable mesh quality around the wing.

To improve resolution in regions with strong flow gradients, two refinement zones were added using refinement boxes:
\begin{itemize}
    \item a \textbf{near-field} refinement box, slightly larger than the delta-wing size, to resolve vortices and strong gradients close to the aircraft,
    \item a \textbf{far-field} refinement box, larger than the near-field zone and extended downstream, to capture the wake behaviour after the flow passes the aircraft.
\end{itemize}
These refinement regions increase mesh density only where needed, keeping the total number of cells at a reasonable level.

\subsubsection{Boundary Layer Mesh}
Boundary layers are one of the most important parts of mesh generation in external aerodynamics, because aerodynamic forces strongly depend on the near-wall velocity gradient. In addition, the layer region often contains a large portion of the final mesh cells.

Two final meshes were generated for comparison:
\begin{itemize}
    \item \textbf{Mesh 1:} target $y^+ \approx 4$,
    \item \textbf{Mesh 2:} target $y^+ \approx 5$.
\end{itemize}
The first cell height was estimated using a $y^+$ calculator, and at least 10 prism layers were added on the wing surface. With 10 layers, the layer growth ratio was kept smaller than 1.15. Growth ratios larger than 1.2 are generally considered poor practice because the transition between layers becomes too aggressive, which can reduce accuracy in the boundary layer and near separation regions.

During the layer generation stage, a practical issue was observed: \texttt{snappyHexMesh} can produce low-quality or collapsed layers when trying to generate more than about 5 layers in a single run. To avoid this, the layer addition was performed in stages:
\begin{enumerate}
    \item \texttt{snappyHexMesh} was first run with \texttt{castellatedMesh} and \texttt{snap} enabled, while \texttt{addLayers} was disabled, producing a clean mesh without layers.
    \item The mesh from the first step was then used as the initial mesh, and \texttt{snappyHexMesh} was run again to add 5 layers.
    \item After verifying layer quality, \texttt{snappyHexMesh} was run once more to add the remaining 5 layers. During this step, the \texttt{finalLayerThickness} parameter was reduced by a factor of two, because half of the total layer thickness already existed from the first layer run.
\end{enumerate}
This staged method produced stable and good-quality layers for both $y^+$ targets.

\subsubsection{Mesh Quality Check}
At the end of the meshing process, the \texttt{checkMesh} utility was used to evaluate mesh quality parameters such as non-orthogonality, skewness, and overall mesh consistency. Both meshes contained a small number of skewed faces, with a maximum skewness value of approximately 11. For large meshes, a small number of locally skewed cells is typically acceptable; however, these regions should be monitored because poor-quality cells can affect solver stability and accuracy.

Finally, the main mesh parameters are summarised in Table~\ref{tab:mesh_parameters} (to be added), including the total number of cells, number of prism layers, first-layer thickness, estimated $y^+$ range for each mesh, and key mesh quality metrics.

% Example placeholder table (fill with your real values):
% \begin{table}[h!]
% \centering
% \caption{Final mesh parameters for both meshes.}
% \label{tab:mesh_parameters}
% \begin{tabular}{lcc}
% \hline
% Parameter & Mesh 1 ($y^+\approx 4$) & Mesh 2 ($y^+\approx 5$) \\
% \hline
% Total cells &  &  \\
% Prism layers & 10 & 10 \\
% Growth ratio & 1.11 & 1.11 \\
% Max skewness & 11 & 11 \\
% \hline
% \end{tabular}
% \end{table}
\section{Boundary Conditions and Solver Settings}
The initial boundary condition files for both turbulence models were taken from the official OpenFOAM \texttt{motorBike} tutorial, which is included in the standard OpenFOAM installation. This tutorial provides a complete and tested setup for incompressible turbulent flow around a solid body using \texttt{simpleFoam}. After copying these files, the boundary conditions were adapted to match the delta wing simulation, including the correct freestream velocity, angle of attack, domain patch names, and turbulence parameter values.

For the velocity and turbulence variable fields at the outer boundaries of the domain (inlet, outlet, upper wall, lower wall, and side wall), the \texttt{freestream} boundary condition type was used. This type automatically switches between inlet and outlet behaviour depending on the local flow direction at each face. At low angles of attack, the flow enters mainly through the front patch and leaves through the rear. However, at high angles of attack, part of the flow may also enter or leave through the upper, lower, or side boundaries. A fixed-velocity condition on these patches would force the flow in an incorrect direction, while a pure zero-gradient condition could cause instability. The \texttt{freestream} type avoids both problems by applying the prescribed freestream value where the flow enters the domain and a zero-gradient condition where it exits. This makes the setup robust across the entire angle-of-attack range without requiring changes for each case.

For the pressure field, the \texttt{freestreamPressure} type was used on the same outer boundaries. This condition adjusts the pressure treatment depending on whether the flow is entering or leaving the domain at each face. At the outlet patch, a fixed reference pressure value is set, which is required for incompressible solvers to define the absolute pressure level. The symmetry plane at the model centreline uses a \texttt{symmetry} boundary condition to simulate only half of the domain.

On the aircraft surface, a no-slip wall condition is applied for velocity. For pressure, a zero-gradient condition is used. For the turbulence variables, wall-function boundary conditions are applied on the aircraft patch. Wall functions allow the first mesh cell to be placed further from the wall than would be needed if the boundary layer were fully resolved. This significantly reduces the total cell count and computational time while still providing a reasonable estimate of the wall shear stress and turbulent quantities. The choice of wall functions must be consistent with the near-wall $y^{+}$ value discussed in Section~3.3.

Two turbulence models were used: the Spalart--Allmaras (SA) model and the $k$--$\omega$ SST model. These were selected because both were also used in the reference ANSYS Fluent simulations by Giannelis et~al.\ \cite{Giannelis2023SSAMGen5}, which allows a direct comparison. Because each model solves different transport equations, the OpenFOAM case requires different files in the initial conditions directory (the \texttt{0/} folder) and the \texttt{system} folder. The SA model requires a single turbulence field file called \texttt{nuTilda}, while the SST model requires two files \texttt{k} and \texttt{omega} instead. All other files, including velocity (\texttt{U}), pressure (\texttt{p}), and turbulent viscosity (\texttt{nut}), remain the same between the two setups. Switching between the models requires replacing these turbulence variable files and changing the model selection in the \texttt{turbulenceProperties} configuration file.

The simulations were run using the \texttt{simpleFoam} solver with second-order accurate numerical schemes for both convective and diffusive terms. Convergence was monitored by tracking residuals and aerodynamic force coefficients. A simulation was considered converged when the residuals dropped below a set threshold and the force coefficients no longer changed between iterations. The same solver settings, numerical schemes, and convergence criteria were kept identical between the SA and SST simulations to ensure that any differences in the results arise only from the turbulence model.

The complete case configuration, including all boundary condition files, solver settings, and numerical schemes, is provided in Annex~[X].

\section{Post-Processing Methods}


For every iteration the aerodynamic force coefficients were extracted using the \texttt{forceCoeffs} function object available in OpenFOAM. The \texttt{forceCoeffs} utility computes the lift coefficient ($C_L$), drag coefficient ($C_D$), and pitching moment coefficient ($C_M$) by integrating the pressure and viscous (shear) contributions over the selected aircraft surface patches. The output is written during the run, which allows the coefficient histories to be monitored directly during the iteration process \cite{OpenFOAMForceCoeffs,OpenFOAMForces}.

To obtain physically meaningful and comparable coefficients, the required reference values were specified in the \texttt{forceCoeffs} setup: freestream velocity magnitude, reference area, reference length (mean aerodynamic chord), and the moment reference point. These values were selected to match the SSAM-Gen5 reference parameters reported by Giannelis et al.~\cite{Giannelis2023SSAMGen5}, so that the resulting coefficients can be compared directly with both the wind tunnel measurements and the ANSYS Fluent results.

The lift and drag directions were defined explicitly in the \texttt{forceCoeffs} configuration. Since the simulations cover a wide range of angles of attack, the lift and drag directions change for each case. Drag was defined parallel to the freestream direction, while lift was defined perpendicular to the freestream direction. Therefore, for each angle of attack, the corresponding lift and drag direction vectors were recalculated and updated before running the case.

Convergence of the aerodynamic coefficients was assessed by plotting the coefficient histories as a function of iteration number. In well-converged cases, $C_L$, $C_D$, and $C_M$ approach steady values and no longer show a systematic drift. For some high angle-of-attack cases, small oscillations remained even after the residuals decreased, which is typical for highly separated flows. In these cases, the final reported coefficients were obtained by averaging over the last several hundred iterations after the oscillations became statistically stable.

Flow visualisation and qualitative inspection were performed using the open-source software ParaView, which can read OpenFOAM case data directly \cite{Ahrens2005ParaView}. ParaView was used mainly to examine pressure (and pressure coefficient) contours on the delta-wing surface and to visualise flow structures in the domain at different angles of attack. These visual checks help confirm that the solution contains physically reasonable features, such as the development of leading-edge vortices at moderate and high angles of attack.

The near-wall resolution was also verified by checking the $y^{+}$ values on the aircraft surface after each simulation. OpenFOAM provides a dedicated \texttt{yPlus} function object that computes and writes the $y^{+}$ field for turbulence models \cite{OpenFOAMyPlus}. The resulting $y^{+}$ distribution was visualised in ParaView to confirm that the values remain consistent with the chosen near-wall treatment and boundary condition setup across the full surface.

All post-processing steps (coefficient extraction, convergence monitoring, flow visualisation, and $y^{+}$ verification) were applied in the same manner for both turbulence models (SA and $k$--$\omega$ SST) and for all angles of attack. This ensures that the comparisons between turbulence models and against the reference data are consistent.



\chapter{Validation Results}\label{chapter:discussion}
\section{Delta Wing Mesh Validation: OpenFOAM vs Wind Tunnel vs ANSYS Fluent}


To verify that the CFD results are not strongly affected by mesh resolution, two meshes with different refinement levels were generated and compared. The first mesh (Mesh~1) contains approximately $8.74\times 10^{6}$ cells with a target $y^{+}\approx 4$, while the second mesh (Mesh~2) contains approximately $1.93\times 10^{6}$ cells with a target $y^{+}\approx 5$. Both meshes were generated with 10 prism layers on the aircraft surface. The mesh comparison was performed using the Spalart--Allmaras (SA) turbulence model, and both meshes were simulated over the full angle of attack range from $-4^{\circ}$ to $30^{\circ}$. The results are summarised in Table~\ref{tab:mesh_SA}.

At low and moderate angles of attack (from $-4^{\circ}$ to approximately $10^{\circ}$), both meshes provide very similar aerodynamic coefficients. In this range, the lift coefficient differs by less than about 1\% between the meshes. For example, at $10^{\circ}$, Mesh~1 predicts $C_{L}=0.5635$ while Mesh~2 predicts $C_{L}=0.5684$, which corresponds to a difference below 0.9\%. The drag coefficient is also nearly identical in this range, with differences typically below 1\%. These results indicate that both meshes provide sufficient resolution when the flow is mostly attached.

At higher angles of attack (above approximately $14^{\circ}$), the differences between the meshes become more visible. The coarser Mesh~2 tends to predict slightly higher lift and drag compared to the finer Mesh~1. At $20^{\circ}$, Mesh~1 predicts $C_{L}=1.0199$ and $C_{D}=0.3822$, while Mesh~2 predicts $C_{L}=1.0722$ and $C_{D}=0.3970$, which corresponds to a difference of roughly 5\% in lift and 4\% in drag. At $30^{\circ}$, the lift values become close again ($C_{L}=1.3211$ for Mesh~1 and $C_{L}=1.3116$ for Mesh~2), and the drag values are also very similar ($C_{D}=0.7599$ for Mesh~1 and $C_{D}=0.7534$ for Mesh~2).

The pitching moment coefficient shows a higher sensitivity to the mesh at high angles of attack, which is expected because the pitching moment depends on the detailed pressure distribution over the full surface and is therefore more affected by local mesh differences. For instance, at $22^{\circ}$, Mesh~1 predicts $C_{My,cog}=0.0626$ while Mesh~2 predicts $C_{My,cog}=0.0820$. Even though the absolute values differ more than for lift and drag, the overall trend of the pitching moment curve remains consistent.

When compared to the wind tunnel data and the ANSYS Fluent results reported by Giannelis et al.~\cite{Giannelis2023SSAMGen5}, both meshes show good agreement in the low to moderate angle of attack range. At high angles of attack, some deviation is expected because steady RANS simulations have known limitations for flows with strong separation, vortex interaction, and possible vortex breakdown. Similar deviations were also reported in the ANSYS Fluent results in the same reference.

Based on these comparisons, Mesh~1 (approximately $8.7$ million cells) was selected for all subsequent simulations. Although Mesh~2 provides reasonable predictions across most of the angle of attack range, the finer mesh gives more consistent results at high angles of attack, especially for the pitching moment coefficient.

\section{Delta Wing Turbulence Model Validation: OpenFOAM vs Wind Tunnel vs ANSYS Fluent}


After selecting the finer mesh (Mesh~1, $8.74\times 10^{6}$ cells, $y^{+}\approx 4$), the influence of the turbulence model was evaluated by comparing the Spalart--Allmaras (SA) and $k$--$\omega$ SST models. The results are summarised in Table~\ref{tab:tm_mesh1}.

At low angles of attack (from $-4^{\circ}$ to approximately $8^{\circ}$), both turbulence models produce very similar results. For example, at $6^{\circ}$, the SA model predicts $C_{L}=0.3375$ and $C_{D}=0.0573$, while the $k$--$\omega$ SST model predicts $C_{L}=0.3284$ and $C_{D}=0.0560$. In this range the flow remains mostly attached, so the solution is less sensitive to the turbulence model choice.

From approximately $10^{\circ}$ and higher, the differences between the models become more visible. In the mid-range between $14^{\circ}$ and $20^{\circ}$, the $k$--$\omega$ SST model predicts slightly higher lift compared to SA (for example, at $20^{\circ}$: $C_{L}=1.0756$ for SST vs $C_{L}=1.0199$ for SA). The drag coefficient follows a similar tendency, with SST predicting slightly higher drag at many high angles.

The largest differences are observed in the pitching moment coefficient at higher angles of attack. This behaviour is expected because the pitching moment is highly sensitive to small changes in pressure distribution, which depend on how each turbulence model represents separation and vortex structures.

Overall, the trends obtained in OpenFOAM are consistent with the behaviour reported by Giannelis et al.~\cite{Giannelis2023SSAMGen5} for ANSYS Fluent, and both turbulence models reproduce the general shape of the lift, drag, and pitching moment curves across the full angle of attack range.

% ------------------------------------------------------------
% Tables
% ------------------------------------------------------------

\begin{table}[!ht]
\centering
\caption{Mesh comparison using the Spalart--Allmaras turbulence model: Mesh~1 ($8.74\times10^{6}$ cells, $y^{+}\approx 4$) vs Mesh~2 ($1.93\times10^{6}$ cells, $y^{+}\approx 5$). Both meshes use 10 prism layers.}
\label{tab:mesh_SA}
\setlength{\tabcolsep}{6pt}
\renewcommand{\arraystretch}{1.15}
\begin{tabular}{c c c c | c c c c}
\hline
\multicolumn{4}{c|}{\textbf{SA --- Mesh 1: $8{,}741{,}586$ cells, $y^{+}\approx 4$, 10 layers}} &
\multicolumn{4}{c}{\textbf{SA --- Mesh 2: $1{,}926{,}908$ cells, $y^{+}\approx 5$, 10 layers}}\\
\hline
$\alpha$ [$^\circ$] & $C_L$ & $C_D$ & $C_{My,cog}$ &
$\alpha$ [$^\circ$] & $C_L$ & $C_D$ & $C_{My,cog}$\\
\hline
-4 & -0.1978 & 0.0362 & -0.0131 & -4 & -0.2343 & 0.0426 & -0.0249 \\
-2 & -0.1181 & 0.0324 & -0.0125 & -2 & -0.1187 & 0.0328 & -0.0123 \\
0  & -0.0091 & 0.0292 & -0.0004 & 0  & -0.0091 & 0.0297 & -0.0006 \\
2  & 0.1007  & 0.0314 & 0.0109  & 2  & 0.1010  & 0.0321 & 0.0100 \\
4  & 0.2144  & 0.0397 & 0.0244  & 4  & 0.2145  & 0.0399 & 0.0232 \\
6  & 0.3375  & 0.0573 & 0.0388  & 6  & 0.3398  & 0.0575 & 0.0365 \\
8  & 0.4570  & 0.0836 & 0.0552  & 8  & 0.4599  & 0.0837 & 0.0529 \\
10 & 0.5635  & 0.1170 & 0.0732  & 10 & 0.5684  & 0.1170 & 0.0729 \\
12 & 0.6743  & 0.1592 & 0.0856  & 12 & 0.6797  & 0.1591 & 0.0834 \\
14 & 0.7901  & 0.2101 & 0.0884  & 14 & 0.8019  & 0.2116 & 0.0854 \\
16 & 0.8828  & 0.2653 & 0.0826  & 16 & 0.9110  & 0.2705 & 0.0843 \\
18 & 0.9679  & 0.3256 & 0.0779  & 18 & 1.0063  & 0.3344 & 0.0862 \\
20 & 1.0199  & 0.3822 & 0.0753  & 20 & 1.0722  & 0.3970 & 0.0899 \\
22 & 1.0853  & 0.4502 & 0.0626  & 22 & 1.1121  & 0.4568 & 0.0820 \\
24 & 1.1309  & 0.5119 & 0.0943  & 24 & 1.1636  & 0.5247 & 0.0781 \\
26 & 1.2046  & 0.5903 & 0.0949  & 26 & 1.2354  & 0.6048 & 0.0784 \\
28 & 1.2489  & 0.6651 & 0.1032  & 28 & 1.2845  & 0.6817 & 0.0822 \\
30 & 1.3211  & 0.7599 & 0.0866  & 30 & 1.3116  & 0.7534 & 0.0844 \\
\hline
\end{tabular}
\end{table}

\begin{table}[!ht]
\centering
\caption{Turbulence model comparison on Mesh~1 ($8.74\times10^{6}$ cells, $y^{+}\approx 4$): Spalart--Allmaras vs $k$--$\omega$ SST.}
\label{tab:tm_mesh1}
\setlength{\tabcolsep}{6pt}
\renewcommand{\arraystretch}{1.15}
\begin{tabular}{c c c c | c c c c}
\hline
\multicolumn{4}{c|}{\textbf{Mesh 1 --- SA}} &
\multicolumn{4}{c}{\textbf{Mesh 1 --- $k$--$\omega$ SST}}\\
\hline
$\alpha$ [$^\circ$] & $C_L$ & $C_D$ & $C_{My,cog}$ &
$\alpha$ [$^\circ$] & $C_L$ & $C_D$ & $C_{My,cog}$\\
\hline
-4 & -0.1978 & 0.0362 & -0.0131 & -4 & -0.2084 & 0.0375 & -0.0081 \\
-2 & -0.1181 & 0.0324 & -0.0125 & -2 & -0.0923 & 0.0292 & -0.0169 \\
0  & -0.0091 & 0.0292 & -0.0004 & 0  & -0.0055 & 0.0277 & 0.0034 \\
2  & 0.1007  & 0.0314 & 0.0109  & 2  & 0.0981  & 0.0281 & 0.0115 \\
4  & 0.2144  & 0.0397 & 0.0244  & 4  & 0.2061  & 0.0382 & 0.0297 \\
6  & 0.3375  & 0.0573 & 0.0388  & 6  & 0.3284  & 0.0560 & 0.0437 \\
8  & 0.4570  & 0.0836 & 0.0552  & 8  & 0.4395  & 0.0814 & 0.0642 \\
10 & 0.5635  & 0.1170 & 0.0732  & 10 & 0.5500  & 0.1149 & 0.0803 \\
12 & 0.6743  & 0.1592 & 0.0856  & 12 & 0.6483  & 0.1536 & 0.0903 \\
14 & 0.7901  & 0.2101 & 0.0884  & 14 & 0.7643  & 0.2053 & 0.0893 \\
16 & 0.8828  & 0.2653 & 0.0826  & 16 & 0.8942  & 0.2667 & 0.0915 \\
18 & 0.9679  & 0.3256 & 0.0779  & 18 & 1.0012  & 0.3320 & 0.0914 \\
20 & 1.0199  & 0.3822 & 0.0753  & 20 & 1.0756  & 0.3980 & 0.1004 \\
22 & 1.0853  & 0.4502 & 0.0626  & 22 & 1.1426  & 0.4684 & 0.1069 \\
24 & 1.1309  & 0.5119 & 0.0943  & 24 & 1.1951  & 0.5366 & 0.0938 \\
26 & 1.2046  & 0.5903 & 0.0949  & 26 & 1.2208  & 0.6002 & 0.0819 \\
28 & 1.2489  & 0.6651 & 0.1032  & 28 & 1.2604  & 0.6720 & 0.0893 \\
30 & 1.3211  & 0.7599 & 0.0866  & 30 & 1.3176  & 0.7634 & 0.0741 \\
\hline
\end{tabular}
\end{table}



\chapter{CAFA Tech Drone Simulation}\label{chapter:discussion}
\section{CAFA Tech Drone Baseline Configuration}
\section{CAFA Tech Drone with External Stores}
\section{Drag Increment Analysis}


\chapter{Summary}\label{chapter:summary} 
\input{chapters/summary.tex}

\zlabel{lastpagetocount}        % DO NOT REMOVE! Used for counting number of pages of main text
